\section{Introduction}
In \textit{Informational Energetics: A Universal Architecture of Persistence} \cite{meyer_informational_2026}, we introduced Informational Energetics (IE) and established that any persistent system must resist entropic decay by satisfying six architectural pillars derived from control theory, information theory, and thermodynamics: Capacity (the vessel), Map (the model), Protocol (the interface), Governor (the constraint), Toll (the temporal cost), and Margin (the buffer). In \textit{Informational Energetics: Entropic Action} \cite{meyer_ieaction_2026}, we formalized the Entropic Action ($S_\Phi$) which defines the total thermodynamic cost for a persistent system to maintain structural coherence.

While previous work established what a persistent system must \textit{be} and how it must \textit{behave}, neither addressed how it can continue to persist when its regulatory complexity exceeds the capacity of a single control loop. Sub-partitioning can be found in every domain: a biological cell that must simultaneously maintain membrane integrity, metabolize, and replicate partitions into organelles; a corporation partitions into departments. In \cref{sec:nested_persistence}, we define three fundamental rules of \textbf{Nested Persistence} governing this partitioning: \textit{Capacity Partitioning}, \textit{Reversible Encoding}, and \textit{Impedance Matching}. Together, they define how a system partitions into specialized subsystems to ensure continued persistence.

\subsection{The Test: The Vacuum}

The vacuum is the ideal test case. Macroscopic systems frequently utilize external energy for partitioning, but a closed system with no external environment to offset its dissipation is forced to solve the partitioning problem with zero error. It therefore provides the most stringent possible test of whether these rules are universal.

To test these rules, we apply them to the $E_8$ substrate derived in \cite{meyer_informational_2026}. The functional operator forms were established in \cite{meyer_ieaction_2026}, but their numerical coefficients remained undetermined. When the rules of Nested Persistence are applied, they emerge as coupled partition coefficients of a single finite capacity budget. Continuing the numbering established in the companion papers (System 0: architectural requirements, System 1: the substrate, System 2: geometric impedance, and System 3: Entropic Lagrangian dynamics):

\begin{itemize}
    \item \textbf{System 4: The Gauge Partition} (Rule 1: Capacity Partitioning). The $\nu=16$ chiral capacity is subdivided into non-overlapping geometric allocations, yielding the gauge couplings: Strong Coupling ($\alpha_s$), Weak Mixing Angle ($\sin^2\theta_W$), and Jarlskog Invariant ($J$). Because the temporal and color sectors are geometrically isolated by this partition, we show that the Strong CP problem is structurally prohibited without requiring an axion.
    
    \item \textbf{System 5: The Resonant Interface} (Rule 3: Impedance Matching). A nested subsystem stabilizing local topological boundaries against vacuum resonance, identified as the Higgs mechanism. Generates the electroweak mass scale: Higgs VEV ($v$), Fermi constant ($G_F$), and Electron Yukawa ($y_e$).
    
    \item \textbf{System 6: The Geometric Depth} (Rules 1 \& 2: Capacity Partitioning and Reversible Encoding). A nested subsystem stabilizing the macroscopic lattice volume, identified as Gravity. Signal attenuation derives the Planck scale ($M_P$) and Cosmological Constant ($\Lambda$), which we show accounts for the observed Gravitational Hierarchy and Vacuum Catastrophe.
\end{itemize}

A zero-degree-of-freedom global geometric fit spanning 64 orders of magnitude ($\alpha^{-1}$, $G_F$, $M_W$, $\alpha_s$, $M_H$, $\sin^2\theta_W$, $J$, $M_P$, $R_K$) yields a reduced $\chi^2 \approx \combinedreducedchiVal$. This convergence demonstrates that the fundamental constants are the structurally necessary partition coefficients of a finite information substrate.